% Chapter 9: Análisis de Riesgo
\chapter{Análisis de Riesgo}

\subsection{Identificación y Mitigación de Riesgos Sistémicos}

La consultoría en ingeniería minera está expuesta tanto a la volatilidad macroeconómica del ciclo de materias primas como a complejidades operacionales específicas de su cartera de proyectos.

\begin{figure}[htbp]
\centering
\begin{tikzpicture}[scale=0.95]
  % Grid cells with fill colors
  \fill[lightgreen!50] (0,0) rectangle (1,1);
  \fill[lightgreen!50] (0,1) rectangle (1,2);
  \fill[lightgreen!50] (1,0) rectangle (2,1);
  
  \fill[lightorange!50] (0,2) rectangle (1,3);
  \fill[lightorange!50] (1,1) rectangle (2,2);
  \fill[lightorange!50] (2,0) rectangle (3,1);
  \fill[lightorange!50] (1,2) rectangle (2,3);
  \fill[lightorange!50] (2,1) rectangle (3,2);
  \fill[lightorange!50] (0,3) rectangle (1,4);
  
  \fill[warningorange!50] (3,0) rectangle (4,1);
  \fill[warningorange!50] (3,1) rectangle (4,2);
  \fill[warningorange!50] (2,2) rectangle (3,3);
  \fill[warningorange!50] (1,3) rectangle (2,4);
  
  \fill[lightred!70] (2,3) rectangle (3,4);
  \fill[lightred!70] (3,2) rectangle (4,3);
  \fill[alertred!60] (3,3) rectangle (4,4);
  
  % Grid lines
  \draw[white, line width=1.5pt] (0,0) grid (4,4);
  
  % Labels
  \draw[line width=1.2pt, ->] (0,0) -- (4.3,0) node[right, font=\small] {Impacto};
  \draw[line width=1.2pt, ->] (0,0) -- (0,4.3) node[above, font=\small] {Probabilidad};
  
  \node[below] at (0.5,0) {\scriptsize Bajo};
  \node[below] at (1.5,0) {\scriptsize Medio};
  \node[below] at (2.5,0) {\scriptsize Alto};
  \node[below] at (3.5,0) {\scriptsize Crítico};
  
  \node[left] at (0,0.5) {\scriptsize Improbable};
  \node[left] at (0,1.5) {\scriptsize Posible};
  \node[left] at (0,2.5) {\scriptsize Probable};
  \node[left] at (0,3.5) {\scriptsize Muy Probable};
  
  % Risks Plotted
  \node[circle, draw=black, fill=white, inner sep=1pt, font=\scriptsize\bfseries] at (2.5,2.5) {R1};
  \node[circle, draw=black, fill=white, inner sep=1pt, font=\scriptsize\bfseries] at (3.5,1.5) {R2};
  \node[circle, draw=black, fill=white, inner sep=1pt, font=\scriptsize\bfseries] at (2.5,2.8) {R3};
  \node[circle, draw=black, fill=white, inner sep=1pt, font=\scriptsize\bfseries] at (1.5,3.5) {R4};
  \node[circle, draw=black, fill=white, inner sep=1pt, font=\scriptsize\bfseries] at (2.5,3.5) {R9};
  
  % Legend on the right
  \node[anchor=west, font=\tiny] at (4.5,3.5) {\textbf{Leyenda:}};
  \node[anchor=west, font=\tiny] at (4.5,3.1) {R1: Volatilidad Precios (Probable, Alto)};
  \node[anchor=west, font=\tiny] at (4.5,2.7) {R2: Sanciones Rusia (Posible, Crítico)};
  \node[anchor=west, font=\tiny] at (4.5,2.3) {R3: Concentración Clientes (Probable, Alto)};
  \node[anchor=west, font=\tiny] at (4.5,1.9) {R4: Escasez de Talento (Muy Probable, Medio)};
  \node[anchor=west, font=\tiny] at (4.5,1.5) {R9: Retrasos en Permitting (Muy Probable, Alto)};
\end{tikzpicture}
\caption{Matriz de Evaluación de Risks para la Consultoría de Ingeniería Minera.}
\label{fig:risk_map}
\end{figure}

A continuación se detallan los 5 riesgos de mayor impacto para REDCO y sus estrategias de mitigación:

\begin{enumerate}
    \item \textbf{R2 - Sanciones Geopolíticas y Concentración en Rusia (Efecto Seligdar):}
    \begin{itemize}
        \item \textbf{Impacto:} Crítico | \textbf{Probabilidad:} Posible.
        \item \textbf{Descripción:} Bloqueo de pagos, imposibilidad de exportar software/servicios de ingeniería a operadores sancionados (Polyus, Seligdar), daño reputacional.
        \item \textbf{Mitigación:} Activación del escenario base ajustado en flujo de caja (ignorar ingresos POM de Rusia para financiamiento operacional). Canalizar cualquier servicio remanente vía intermediarios en Asia Central o China que actúen como buffer legal.
    \end{itemize}

    \item \textbf{R3 - Concentración de Clientes:}
    \begin{itemize}
        \item \textbf{Impacto:} Alto | \textbf{Probabilidad:} Likely.
        \item \textbf{Descripción:} Dependencia de 1 o 2 grandes proyectos o clientes (p.ej., REDCO dependiendo excesivamente del POM ruso o de un único operador en Chile).
        \item \textbf{Mitigación:} Diversificación comercial agresiva hacia Perú y Brasil. Establecer un límite máximo de exposición de facturación por cliente del 25\%.
    \end{itemize}

    \item \textbf{R4 - Escasez de Talento Clave:}
    \begin{itemize}
        \item \textbf{Impacto:} Medio | \textbf{Probabilidad:} Very Likely.
        \item \textbf{Descripción:} Dificultad para reclutar y retener ingenieros de procesos, planificadores mineros y geotecnistas de alto nivel frente a las ofertas de operadores mineros directos.
        \item \textbf{Mitigación:} Desarrollar el baseline organizacional de REDCO (Pilar B - Capability × Personal Identity) para estructurar planes de carrera, flexibilidad laboral y bonos por desempeño alineados al margen de proyectos.
    \end{itemize}

    \item \textbf{R9 - Retrasos en Permitting de Proyectos de Clientes:}
    \begin{itemize}
        \item \textbf{Impacto:} Alto | \textbf{Probabilidad:} Very Likely.
        \item \textbf{Descripción:} Los clientes congelan la ingeniería de detalle porque sus EIAs están atrapados en el SEA (Chile) o Senace (Perú).
        \item \textbf{Mitigación:} Desarrollar capacidades internas de pre-permitting e ingeniería regulatoria para integrarlas en etapas tempranas y ayudar al cliente a acelerar la obtención de permisos.
    \end{itemize}

    \item \textbf{R1 - Volatilidad del Precio de los Metales:}
    \begin{itemize}
        \item \textbf{Impacto:} Alto | \textbf{Probabilidad:} Likely.
        \item \textbf{Descripción:} Caída drástica del cobre o hierro reduce el CAPEX de los clientes, cancelando proyectos de ingeniería.
        \item \textbf{Mitigación:} Pivotar la propuesta de valor en periodos de precios bajos hacia OPEX (ingeniería de optimización de costos, eficiencia hídrica y de procesos operacionales) que las mineras contratan activamente para mantener márgenes.
    \end{itemize}
\end{enumerate}
